\chapter{Conclusion}\label{ch:conclusion}

\section{Summary of Findings}

This thesis investigated streaming communication for scalable, policy-aware data exchange systems, using DYNAMOS as a case study. The central problem was not simply that large messages are slow. The deeper issue is that policy-aware middleware must improve bulk transfer without weakening the orchestration path that decides which exchange is allowed, where it may run, and which components participate.

RQ1 asked which communication mechanisms are suitable for large-scale and continuous data exchange in microservice-based systems. The controlled benchmark shows that direct gRPC transfer is strongest for local service-to-service communication, with client-streaming giving the highest clean throughput. Batched unary is close in the clean matrix and stronger in the backpressure scenario. Brokered transports are slower locally, but they provide decoupling, retention, replay, and failure-handling options that direct streams do not provide by themselves.

RQ2 asked how streaming communication can be integrated while preserving dynamic service composition, policy-driven orchestration, and isolation. The DYNAMOS implementation keeps the existing control path intact: requests are still evaluated and composed before execution, and generated jobs remain the execution unit. The change is in the approved bulk data path. Provider-side chunks must remain chunks through aggregate, algorithm, agent, and gateway stages; otherwise the system merely hides a full-buffered result behind a streaming interface.

RQ3 validated the implemented variants on local bulk \texttt{dataThroughTtp} retrieval. All chunk-aware variants beat the classic unary baseline in average completion time and peak memory for 50k, 100k, and 250k row result sets. They also expose earlier partial results, which classic unary cannot do because its first result is its final result. Focused checks on \texttt{computeToData} and multi-provider \texttt{dataThroughTtp} show that the implementation is not confined to one request shape. The strongest local recommendation is chunked unary. It gives the best 100k and 250k averages, the lowest 250k memory peak, and the smallest architectural shift from the original message-driven design.

The overall answer is therefore boundary-specific. Streaming should be treated as a data-path change after policy approval, not as a replacement for policy enforcement or orchestration. Direct gRPC streaming is useful where generated services have naturally coupled lifetimes. Chunked unary is the best current local default for DYNAMOS because it preserves the existing request-shaped architecture while removing full-result buffering. RabbitMQ Streams remains a distributed-boundary candidate, where retained records, replay, and independent consumers may matter more than local latency.

\section{Contributions}

This thesis makes four main contributions. First, it provides a controlled comparison of six transport mappings for large-scale and continuous microservice data exchange: client-streaming gRPC, per-message unary gRPC, batched-unary gRPC, RabbitMQ Streams, NATS JetStream, and Kafka. Second, it identifies chunk preservation as the central architectural requirement for improving large-result transfer in policy-aware generated workflows. Third, it implements and evaluates three chunk-aware DYNAMOS data paths: chunked unary transfer, intra-job gRPC streaming, and RabbitMQ Streams at the agent boundary. Fourth, it provides a benchmark corpus and appendices that document the relevant run scopes, parameter rationale, environment, provenance, failure modes, clean-matrix results, and resource measurements.

\section{Limitations}

The conclusions are strongest for the local Kubernetes environments and workloads evaluated in this thesis. The RQ1 benchmark isolates transport overhead on a single node, while the RQ3 benchmark validates DYNAMOS bulk retrieval locally. Cross-site latency, broker placement, replication, independent administrative domains, broader DYNAMOS archetypes, and mid-stream policy changes remain outside the validated scope. Resource usage is sampled rather than continuously traced, and the selected 5{,}000-row chunk size is an experimental configuration rather than a global optimum. These limitations do not invalidate the observed local improvements, but they bound the generality of the recommendations.

\section{Future Work}

The most important future work is distributed validation. The local DYNAMOS benchmark proves that the basic chunk-aware path works before it is meaningful to test it across sites. A natural next step is Scattered Directive, which extends DYNAMOS with cross-node microservice chains for policy-aware VFL workflows and uses FABRIC to automate distributed deployment \cite{jongejans2025scattered,scatteredDirectiveRepo2026}.

In that setting, bulk transfer can become part of a larger VFL-style workflow. Parties may hold different features for overlapping entities, and a TTP or shared execution component may coordinate aligned intermediate values, labels, embeddings, model updates, or other training artefacts. The current thesis does not claim to validate a complete VFL training workload. It validates the communication path that such a workload would depend on when large intermediate results have to cross generated services and agent boundaries.

FABRIC provides programmable compute and storage resources connected by high-speed dedicated links, which makes it a suitable environment for testing policy-aware data exchange when the TTP and data providers are separated across sites \cite{baldin2019fabric,fabricPortal2026}. A later FABRIC experiment should keep the RQ3 metrics from this section and add network throughput, broker lag or stream backlog, cross-site tail latency, failure recovery, concurrent-job scaling, and policy-change responsiveness. RabbitMQ Streams deserves special attention in that setting because its local overhead may be outweighed by retained-stream and replay semantics once agents are no longer co-located.

The second future-work direction is stronger policy handling during active streams. The evaluated implementation preserves pre-stream policy approval, but it does not solve mid-stream revocation or continuous policy enforcement. The usage-control literature provides the target model: ongoing authorisation with revocation during use~\cite{park2004ucon,hariri2023uconplus}. A concrete next step is chunk-boundary re-validation as outlined in \Cref{sec:policy-implications}, followed by lease- or capability-based transfer rights in which the policy engine issues a time- or volume-bounded transfer capability that sidecars enforce and that expires unless renewed. A complete policy-aware streaming middleware should further define how policy changes interrupt active generated jobs, how already delivered chunks are audited, and how downstream services should react when a long-running approved exchange becomes invalid. A third, smaller direction is backpressure-aware orchestration: stream lag and backlog signals from the data path could inform orchestration decisions such as admission and placement, connecting the streaming layer back to the adaptive part of the middleware.
